\section{Gamma-weighted integrals and explicit polygamma bridges}
\label{rlc:sec:gamma}

When $\Rea a>1/2$, the symmetric vertical line $p=1/2+it$ in
\eqref{rlc:eq:barnes} is admissible. Put $b=a-1/2$. Euler reflection
then gives
\begin{equation}\label{rlc:eq:gammacentral}
 \boxed{\displaystyle
 \begin{aligned}
 \CC_{r,a}(W;0)&=\frac1{2\pi}\int_{-\infty}^{\infty}
 \left[\Gamma\left(\tfrac12+it\right)
       \Gamma\left(\tfrac12-it\right)\right]^r
 (b+it)^{-W}\dd t\\
 &=\frac1{2\pi}\int_\R
      \left(\frac\pi{\cosh\pi t}\right)^r(b+it)^{-W}\dd t.
 \end{aligned}}
\end{equation}
The logarithm of $b+it$ is the right-half-plane branch. For
$0<\Rea a\le1/2$ use a nonsymmetric line satisfying
\eqref{rlc:eq:cstrip}; simply keeping the symmetric line can cross
$\Rea(p+a-1)=0$ and is not the stated formula.

Theorem~\ref{rlc:thm:central} now evaluates every integral in
\eqref{rlc:eq:gammacentral} by a finite Hurwitz-zeta expression.
For $r=2$ and $M\ge1$, differentiation in $W$ at zero gives
\begin{equation}\label{rlc:eq:gammastieltjes}
 \boxed{\displaystyle
 \frac\pi2\int_\R
     \frac{\log^M(a-\tfrac12+it)}{\cosh^2\pi t}\dd t
       =-M\gamma_{M-1}(a).}
\end{equation}
There is no alternating sign on the right: the $(-1)^M$ from the
power derivative cancels all but a minus sign in
\eqref{rlc:eq:Ejets}. This is a consistency bridge to classical
Stieltjes integral representations, not a claim of first discovery
of such representations.

\subsection{An ordinary log-Gamma identity}
At $W=-1$, differentiation of
$\E_a(W)=W\zeta(W+1,a)$ yields
\[
 \E_a'(-1)=\zeta(0,a)-\zeta^{[1]}(0,a).
\]
Using \eqref{rlc:eq:gammacentral} and
$\zeta^{[1]}(0,a)=\log\Gamma(a)-\tfrac12\log(2\pi)$ gives
\begin{equation}\label{rlc:eq:gammalog}
 \frac\pi2\int_\R\frac{
 (b+it)\log(b+it)-(b+it)}{\cosh^2\pi t}\dd t
      =\log\Gamma(a)-\frac12\log(2\pi).
\end{equation}
All differentiations are dominated by the hyperbolic exponential
weight. The classical Hurwitz--Gamma identity is the only additional
input \cite{DLMF}.

\subsection{Higher-depth Gamma and polygamma identities}
The all-depth jets in Section~\ref{rlc:sec:stieltjes} give, for example,
\begin{align}
 \frac1{2\pi}\int_\R
 \left(\frac\pi{\cosh\pi t}\right)^4\log(b+it)\dd t
 &=\frac{2\pi^2}{3}\psi(a)+\frac16\psi^{(2)}(a),
 \label{rlc:eq:gamma4log}\\
 \frac1{2\pi}\int_\R
 \left(\frac\pi{\cosh\pi t}\right)^4\log^2(b+it)\dd t
 &=-\frac{4\pi^2}{3}\gamma_1(a)
     +\zeta(3,a)+\frac23\zeta^{[1]}(3,a).
 \label{rlc:eq:gamma4log2}
\end{align}
The first uses $\psi^{(2)}(a)=-2\zeta(3,a)$.
Positive integer total orders likewise reduce to polygamma values by
\[
 \zeta(k,a)=\frac{(-1)^k}{(k-1)!}\psi^{(k-1)}(a),\qquad k\ge2.
\]
At integer shifts these combine with
$\zeta(k,N+1)=\zeta(k)-H_N^{(k)}$; at rational shifts they give finite
root-of-unity coordinates after the usual residue-class filter.

One can also evaluate every polynomial $t$ moment in
\eqref{rlc:eq:gammacentral}. Expand
$it=(b+it)-b$ to obtain
\begin{equation}\label{rlc:eq:tGamma}
 \frac1{2\pi}\int_\R
 \left(\frac\pi{\cosh\pi t}\right)^r(it)^q(b+it)^{-W}\dd t
 =\sum_{j=0}^{q}\binom qj(-b)^{q-j}\CC_{r,a}(W-j;0).
\end{equation}
Subsequent order derivatives insert arbitrary powers of $\log(b+it)$.
Thus the same finite central-factorial family controls Gamma weights,
polynomial moments, polygamma values, and Stieltjes order derivatives.
